%!TEX root = ../../note.tex
\subsection{Bounded Sets; Supremum and Infimum}\label{sec:bounds}

Bounds make sense in any ordered field. What is special about $\R$ is that
a set bounded above always has a \emph{least} upper bound in $\R$; this
section introduces the language needed to say so.

\begin{definition}\label{def:bounded}
    A nonempty set $S \subseteq \R$ is
    \begin{enumerate}[label=(\arabic*)]
        \item \term{bounded above} if there is $u \in \R$ (an \emph{upper
        bound}) with $x \leq u$ for all $x \in S$;
        \item \term{bounded below} if there is $w \in \R$ (a \emph{lower
        bound}) with $w \leq x$ for all $x \in S$;
        \item \term{bounded} if it is bounded above and below, and
        \emph{unbounded} otherwise.
    \end{enumerate}
\end{definition}

\begin{definition}\label{def:sup}
    An upper bound $u_0$ of $S$ is called the \term{supremum} (least upper
    bound) of $S$, written $u_0 = \sup S$, if $u_0 \leq u$ for every upper
    bound $u$ of $S$.

    A lower bound $w_0$ of $S$ is called the \term{infimum} (greatest lower
    bound) of $S$, written $w_0 = \inf S$, if $w_0 \geq w$ for every lower
    bound $w$ of $S$.
\end{definition}

The notation $\sup S$ is justified because a set has at most one supremum:
if $u_1,u_2$ are both suprema, then $u_1\leq u_2$ (as $u_2$ is an upper
bound and $u_1$ is least) and likewise $u_2\leq u_1$. The same holds for
the infimum.

Checking that $u$ is the \emph{least} upper bound requires comparing $u$ with
every other upper bound. The following lemma replaces this by a condition on
$S$ itself: no matter how small a margin $\epsilon$ we allow, $S$ has a point
within $\epsilon$ of $u$.

\begin{lemma}[$\epsilon$-characterization of the supremum]\label{lem:eps-sup}
    An upper bound $u$ of a nonempty set $S\subseteq\R$ is the supremum of
    $S$ if and only if
    \begin{equation*}
        (\forall \epsilon > 0)(\exists x_\epsilon \in S)\,[u - \epsilon < x_\epsilon].
    \end{equation*}
\end{lemma}

\begin{proof}
    ($\Rightarrow$) Let $u=\sup S$ and $\epsilon>0$. Since
    $u-\epsilon<u$ and $u$ is the least upper bound, $u-\epsilon$ is not an
    upper bound, so some $x_\epsilon\in S$ satisfies $u-\epsilon<x_\epsilon$.

    ($\Leftarrow$) If $u$ were not least, there would be an upper bound
    $w<u$. For $\epsilon=u-w>0$ we get $x_\epsilon\in S$ with
    $w=u-\epsilon<x_\epsilon$, so $w$ is not an upper bound after all.
\end{proof}

Thus, to prove $u=\sup S$ one checks two things:
(1) $u$ is an upper bound of $S$;
(2) for every $\epsilon>0$ some $x\in S$ satisfies $u-\epsilon<x\leq u$.
For $\inf S$ reverse the inequalities.

\begin{example}\label{ex:half-open}
    Let $S = \set{x \in \R : 0 \leq x < 1}$. Then $\sup S = 1$ and
    $\inf S = 0$.
\end{example}

\begin{proof}
    \emph{$\sup S=1$.} Clearly $1$ is an upper bound. By completeness
    $u=\sup S$ exists, and $u\leq1$. Suppose $u<1$ and put
    $\epsilon=\frac{1-u}{2}>0$. Then
    \begin{equation*}
        u<u+\epsilon=\frac{u+1}{2}<1,
    \end{equation*}
    and $u+\epsilon\geq0$ because $u\geq0$ (as $0\in S$). So
    $u+\epsilon\in S$ exceeds the upper bound $u$, a contradiction. Hence
    $u=1$.

    \emph{$\inf S=0$.} Clearly $0$ is a lower bound. Every lower bound $w$
    satisfies $w\leq0$ because $0\in S$. Hence $0$ is the greatest lower
    bound.
\end{proof}

\begin{remark}
    In this example $\sup S=1\notin S$, so $S$ has no largest element,
    whereas $\inf S=0\in S$ is also its least element. A supremum is a
    \emph{maximum} exactly when it belongs to the set. Note also that
    ``bounded'' is a property of the set $S$, while $\sup S$ and $\inf S$
    are numbers attached to it.
\end{remark}
